% Section introducing the AgriGov Market platform
\section{Introduction}

While the strategic importance of digitalizing Algeria's agricultural sector is well-recognized, the transition from traditional, informal market structures to a transparent digital ecosystem requires a rigorous understanding of the field's functional needs. The current reliance on over-intermediation and asymmetric price information \cite{CIHEAM2025} creates a complex environment that the \textbf{AgriGov Market} platform aims to restructure. To build a system that effectively bridges the gap between farmers, transporters, and the Ministry of Agriculture, it is first necessary to conduct a thorough analysis of the requirements and the existing operational challenges.

The objective of this chapter is to lay the analytical foundation for the proposed system. By exploring the requirements from both a technical and a business perspective, we identify the specific boundaries of the system and define how various actors will interact within this new digital framework. This phase is crucial for ensuring that the final software solution is not only technically sound but also aligned with the strategic objectives of the Ministry of Agriculture.

The remainder of this chapter is organized as follows. First, we provide a detailed study of existing systems and their limitations to justify our proposed approach. We then proceed to identify the various actors involved and present a static context diagram to visualize the system's environment. Subsequently, we identify and model the core use cases, providing detailed textual specifications and system sequence diagrams for the most critical functionalities. The chapter also previews the primary human-machine interfaces before concluding with a synthesis of the requirements that will guide the design phase.